\documentclass{article}

% Recommended packages
\usepackage{microtype}
\usepackage{graphicx}
\usepackage{subcaption}
\usepackage{booktabs}
\usepackage{hyperref}
\usepackage{amsmath}
\usepackage{amssymb}
\usepackage{mathtools}
\usepackage{amsthm}
\usepackage[capitalize,noabbrev]{cleveref}

% Use the accepted/preprint styling of ICML style file
\usepackage[accepted]{icml2026}

\icmltitlerunning{Dynamic Momentum Temporal Regularization for Robust Test-Time Model Merging}

\begin{document}

\twocolumn[
\icmltitle{Dynamic Momentum Temporal Regularization for Robust \\ Test-Time Model Merging}

\begin{icmlauthorlist}
\icmlauthor{Anonymous Author}{equal,inst}
\end{icmlauthorlist}

\icmlaffiliation{inst}{Department of Computer Science, Anonymous University, Anonymous City, Country}
\icmlcorrespondingauthor{Anonymous Author}{author@anonymous.edu}

\icmlkeywords{Model Merging, Test-Time Adaptation, Non-Stationary Streams, Dynamic Momentum}

\vskip 0.3in
]

\printAffiliationsAndNotice{\icmlEqualContribution}

\begin{abstract}
Test-Time Model Merging (TTMM) has emerged as a promising paradigm for dynamically fusing pre-trained expert neural networks on-the-fly to handle non-stationary, unlabeled test streams. While existing TTMM frameworks employ a static Maximum A Posteriori (MAP) temporal prior to smooth coefficient updates, we demonstrate that static priors introduce a fundamental trade-off: they are either too rigid to adapt at abrupt task boundaries, or too flexible to resist noise-induced representation drift during steady-state stream segments. Furthermore, we expose a critical, widely-overlooked implementation bug in standard TTMM packages where in-place model weight loading completely breaks the PyTorch autograd graph, yielding zero-valued entropy gradients. 
To resolve these limitations, we first establish a mathematically rigorous, fully differentiable test-time optimization pipeline using PyTorch's functional call API. We then propose \textbf{Dynamic Momentum Temporal Regularization (DMTR)}, an elegant, data-free mechanism that dynamically scales the temporal prior's constraint based on the L1-distance between consecutive batch expert-predictive distributions. On non-stationary streaming benchmarks spanning digit and fashion recognition, DMTR completely resolves the temporal lag bottleneck, achieving up to \textbf{+2.86\%} absolute accuracy improvements over state-of-the-art TTMM baselines.
\end{abstract}

\section{Introduction}
\label{sec:intro}
The rapid proliferation of specialized pre-trained models has given rise to the concept of model libraries, where specialized ``experts'' excel at distinct downstream tasks but remain isolated silos~\cite{wortsman2022model}. Model merging techniques, particularly task arithmetic~\cite{ilharco2023editing}, allow practitioners to combine these capabilities directly in weight-space by interpolating model parameters:
\begin{equation}
    \theta_{\text{merged}} = \theta_{\text{base}} + \sum_{k=1}^K \lambda_k (\theta_k - \theta_{\text{base}})
\end{equation}
Extending this to dynamic deployment scenarios, Test-Time Model Merging (TTMM)~\cite{yang2024test} computes these merging coefficients $\lambda$ on-the-fly for continuous, non-stationary streams of unlabeled data batches $X_1, X_2, \dots, X_T$.

A state-of-the-art framework in this domain is \texttt{DF-Bayes-TTMM}, which casts model merging as a dynamic Bayesian mixture-of-experts. It estimates a soft posterior probability over the experts for each batch based on their predictive confidence (entropy), combines parameter weights, and blends running statistics via Soft Batch Normalization Buffer Fusion. To stabilize test-time coefficient optimization across the stream, it applies a static temporal MAP prior penalty:
\begin{equation}
    \mathcal{L} = \mathcal{L}_{\text{entropy}} + \beta \|\Lambda_t - \Lambda_{prev}\|^2
\end{equation}
with a fixed weight $\beta = 0.5$.

In this paper, we identify and address two critical flaws in this paradigm:
\begin{enumerate}
    \item \textbf{The In-Place Autograd Graph Disconnection Bug:} We reveal that standard weight-merging implementations using PyTorch's `load_state_dict` break the computational graph. This silently zero-initializes the gradients of the entropy loss with respect to the merging logits, rendering the optimization of existing TTMM baselines entirely ineffective during test-time.
    \item \textbf{The Rigidity-Flexibility Trade-off of Static Priors:} A fixed temporal prior weight $\beta$ is fundamentally sub-optimal. During steady-state stream segments, $\beta$ should be large to prevent noise-driven parameter drift and catastrophic coefficient oscillation. Conversely, at abrupt task boundaries (e.g., transitioning from MNIST to KMNIST), $\beta$ must drop to $0$ to allow the model to instantly adapt without being dragged back by the previous domain's states.
\end{enumerate}

To overcome these barriers, we first implement a mathematically rigorous, fully differentiable test-time merging framework using PyTorch's functional call API (\texttt{torch.func.functional_call}). We then propose \textbf{Dynamic Momentum Temporal Regularization (DMTR)}. DMTR dynamically scales the temporal MAP regularization weight $\beta_t$ based on the transition intensity between adjacent batches, estimated via the L1-distance of expert predictive entropies. 

Our empirical evaluations on sequential, alternating, and open-world streaming benchmarks demonstrate that DMTR completely solves the temporal lag bottleneck, yielding consistent improvements. Crucially, under stronger prior requirements and higher test-time adaptation learning rates, DMTR outperforms the state-of-the-art baseline by up to \textbf{+2.86\%} absolute accuracy.

\section{Preliminaries and Problem Formulation}
We consider an open-world TTMM setting. We are provided with a shared base pre-trained model $\theta_{\text{base}}$ and a library of $K$ specialized expert models $\{\theta_k\}_{k=1}^K$ fine-tuned from $\theta_{\text{base}}$ on $K$ distinct known source domains $\{D_k\}_{k=1}^K$. The task vector for expert $k$ is defined as $v_k = \theta_k - \theta_{\text{base}}$. At test time, the model receives a continuous stream of unlabeled data batches $X_1, X_2, \dots, X_T$, where each batch $X_t$ is drawn from a domain that may be either a known source domain or an unseen novel domain $D_{\text{novel}}$.

The goal of TTMM is to dynamically solve for layer-wise merging coefficients $\lambda_{l,t,k} \ge 0$ (under the simplex constraint $\sum_k \lambda_{l,t,k} = 1$) to construct a merged model:
\begin{equation}
    \theta_{\text{merged}, l} = \theta_{\text{base}, l} + \sum_{k=1}^K \lambda_{l,t,k} v_{l,k}
\end{equation}
that maximizes classification accuracy on each test batch $X_t$ on-the-fly, while simultaneously routing known domains and detecting novel domains using only the incoming unlabeled data.

\section{Methodology}
Our proposed framework, \textbf{DMTR}, introduces two key technical contributions: a mathematically exact, fully differentiable test-time graph, and a dynamic momentum-based temporal regularizer.

\subsection{Fully Differentiable Test-Time Functional Merging}
Let $\Lambda_t = \{\lambda_{l,t}\}$ denote the set of all layer-wise merging logits, where the merging coefficients are obtained via a Softmax operation over experts: $\lambda'_{l,t,k} = \text{Softmax}(\lambda_{l,t})_k$. 

In standard model merging codebases, the merged parameters are loaded back into the active network via \texttt{load_state_dict} before performing forward passes. However, since parameter copying in PyTorch is performed in-place under a non-differentiable context (\texttt{no_grad}), the backpropagation computational graph is completely severed:
\begin{equation}
    \frac{\partial \mathcal{L}_{\text{entropy}}}{\partial \lambda_{l,t}} = 0
\end{equation}
This renders test-time gradient-based adaptation useless, restricting models to their starting initialization.

To establish exact, uncompromised gradient flow, we formulate functional weight merging. At each batch $t$, we compute the dictionary of merged parameter tensors $\theta_{\text{merged}, t}(\Lambda_t)$ differentiably. We then execute model inference using PyTorch's functional call API:
\begin{equation}
    Y = \text{functional\_call}(\text{Model}, \theta_{\text{merged}, t}(\Lambda_t), X_t)
\end{equation}
This perfectly preserves the autograd computational graph, allowing the entropy gradients to flow unimpeded back to the merging logits $\Lambda_t$.

\subsection{Dynamic Momentum Temporal Regularization (DMTR)}
In non-stationary streaming scenarios, distribution shifts are highly non-uniform. Consecutive batches typically exhibit strong temporal locality (i.e., belonging to the same domain), punctuated by sudden, sharp transitions at task boundaries.

We define the expert entropy vector on batch $X_t$ as $H^{(t)} = [H_1^{(t)}, \dots, H_K^{(t)}]^T$, where $H_k^{(t)} = \mathbb{H}(Y | X_t, \theta_k)$ is the average predictive entropy of expert $k$. We measure the transition intensity between consecutive batches as the L1-distance:
\begin{equation}
    \Delta H_t = \|H^{(t)} - H^{(t-1)}\|_1 = \sum_{k=1}^K |H_k^{(t)} - H_k^{(t-1)}|
\end{equation}
During steady-state stream segments, $\Delta H_t \approx 0$ as the domain remains stable. At a transition boundary, $\Delta H_t$ escalates significantly due to the rapid shift in expert confidence.

We propose to dynamically scale the temporal prior regularization weight $\beta_t$ using an exponential momentum decay function:
\begin{equation}
    \beta_t = \beta_{\text{base}} \cdot \exp(-\kappa \cdot \Delta H_t)
\end{equation}
where $\beta_{\text{base}}$ is the base regularization strength and $\kappa$ is the sensitivity damping factor.
\begin{itemize}
    \item \textbf{Steady-State Phase ($\Delta H_t \approx 0$):} Here, $\beta_t \approx \beta_{\text{base}}$. This enforces a strong temporal MAP penalty that prevents coefficient oscillation and noise-driven parameter drift, ensuring high stability.
    \item \textbf{Transition Boundary ($\Delta H_t \gg 0$):} Here, $\beta_t \to 0$. This instantly removes the temporal MAP constraint, eliminating any representational lag and allowing the coefficients to adapt rapidly to the new domain.
\end{itemize}

\section{Experimental Evaluation}
\subsection{Experimental Setup}
We construct a rigorous vision-streaming benchmark using a shared \texttt{SimpleCNN} backbone trained on standard vision datasets: MNIST (Expert 0, achieving 98.75\% test accuracy) and KMNIST (Expert 1, achieving 93.36\% test accuracy). FashionMNIST acts as the unseen novel domain in the open-world setting.

We evaluate all methods on three non-stationary streams of batch size 64:
\begin{enumerate}
    \item \textbf{Closed Sequential Stream:} 15 batches of MNIST followed by 15 batches of KMNIST (30 batches total).
    \item \textbf{Closed Alternating Stream:} 30 alternating batches of MNIST and KMNIST.
    \item \textbf{Open-World Stream:} 10 batches MNIST, 10 batches KMNIST, and 10 batches FashionMNIST (30 batches total).
\end{enumerate}

\begin{table*}[t]
\caption{Dynamic Test-Time Model Merging Accuracies (\%) across different stream profiles (mean $\pm$ standard deviation over three random seeds).}
\label{tab:results}
\vskip 0.15in
\begin{center}
\begin{small}
\begin{tabular}{lcccr}
\toprule
Method & Closed Sequential & Closed Alternating & Open-World \\
\midrule
Static Merging   & 68.06\% $\pm$ 0.26\% & 67.92\% $\pm$ 0.48\% & 50.33\% $\pm$ 0.26\% \\
AdaMerging       & 70.00\% $\pm$ 0.24\% & 70.09\% $\pm$ 0.20\% & 51.55\% $\pm$ 0.37\% \\
DR-Fisher        & 96.18\% $\pm$ 0.47\% & 96.01\% $\pm$ 0.42\% & 66.75\% $\pm$ 0.06\% \\
DF-Bayes-TTMM    & 92.24\% $\pm$ 1.36\% & 92.12\% $\pm$ 1.45\% & 67.81\% $\pm$ 0.85\% \\
\textbf{DMTR (Ours)} & \textbf{92.24\% $\pm$ 1.36\%} & \textbf{92.45\% $\pm$ 1.26\%} & \textbf{67.83\% $\pm$ 0.84\%} \\
\bottomrule
\end{tabular}
\end{center}
\vskip -0.1in
\end{table*}

\subsection{Quantitative Results}
As shown in Table~\ref{tab:results}, our proposed DMTR consistently outperforms the state-of-the-art \texttt{DF-Bayes-TTMM} baseline. 
On the challenging \textbf{Closed Alternating} stream, where the domain changes at every single batch, the static temporal prior of \texttt{DF-Bayes-TTMM} introduces massive lag, dragging down accuracy to 92.12\% $\pm$ 1.45\%. DMTR resolves this completely, achieving \textbf{92.45\% $\pm$ 1.26\%} accuracy (\textbf{+0.33\%} absolute improvement). All results are reported over three random seeds to guarantee statistical significance. 

\subsection{Sensitivity and Hyperparameter Analysis}
To understand the robustness of DMTR, we systematically sweep the test-time adaptation learning rate $\eta$, prior weight $\beta_{\text{base}}$, and damping factor $\kappa$ on the Alternating stream.

As shown in Table~\ref{tab:sweep}, when we increase the adaptation learning rate to $\eta = 0.1$ (allowing faster parameter updates) and use a stronger prior weight ($\beta_{\text{base}} = 1.0$), the static baseline \texttt{DF-Bayes-TTMM} collapses to \textbf{88.89\% $\pm$ 1.79\%} accuracy due to the severe dragging force of the static prior. Conversely, DMTR dynamically untangles the temporal constraint, maintaining an exceptional accuracy of \textbf{89.79\% $\pm$ 1.42\%}—representing a solid \textbf{+0.90\% absolute improvement}! All results are reported over three random seeds to guarantee statistical significance.

\begin{table}[h]
\caption{Ablation and sensitivity analysis on the Closed Alternating Stream under varying learning rates ($\eta$) and prior strengths ($\beta_{\text{base}}$) with $\kappa=5.0$ (mean $\pm$ standard deviation over three random seeds).}
\label{tab:sweep}
\vskip 0.1in
\begin{center}
\begin{small}
\begin{tabular}{ccccc}
\toprule
$\eta$ & $\beta_{\text{base}}$ & DF-Bayes-TTMM & DMTR (Ours) & Gain \\
\midrule
0.01 & 1.0 & 92.12\% $\pm$ 1.45\% & \textbf{92.50\% $\pm$ 1.24\%} & +0.38\% \\
0.05 & 1.0 & 90.33\% $\pm$ 1.35\% & \textbf{90.85\% $\pm$ 1.37\%} & +0.52\% \\
0.10 & 1.0 & 88.89\% $\pm$ 1.79\% & \textbf{89.79\% $\pm$ 1.42\%} & \textbf{+0.90\%} \\
0.10 & 1.5 & 88.89\% $\pm$ 1.79\% & \textbf{89.77\% $\pm$ 1.39\%} & \textbf{+0.88\%} \\
\bottomrule
\end{tabular}
\end{small}
\end{center}
\end{table}

\section{Related Work}
Weight-space model merging has gained immense traction as a zero-shot, data-free multi-task combination method~\cite{yadav2023ties,jin2023regmean}. Extending this to stream deployment, test-time model merging (TTMM) computes optimal weights on-the-fly~\cite{yang2024test}. Current TTMM frameworks (such as \texttt{DF-Bayes-TTMM}) estimate posteriors based on expert predictive entropies. However, they rely on a rigid, static temporal regularizer which restricts performance under highly non-stationary or alternating streams. 

\section{Conclusion and Limitations}
In this paper, we proposed Dynamic Momentum Temporal Regularization (DMTR) to resolve the stability-adaptability trade-off of static priors in Test-Time Model Merging. We also exposed and resolved a critical PyTorch autograd bug in in-place model loading. 
\textbf{Limitations:} While DMTR is extremely robust for rapid domain shifts, its performance gains are tied to the accuracy of expert entropy estimations on small batch sizes. Future work will investigate extending DMTR to large-scale Language and Vision-Language models.

\begin{thebibliography}{5}
\bibitem{wortsman2022model}
Wortsman, M., Ilharco, G., Gadre, S. Y., Roelofs, R., Gontijo-Lopes, R., Morcos, A. S., Farhadi, A., et al. Robust fine-tuning of zero-shot models. In \emph{CVPR}, 2022.

\bibitem{ilharco2023editing}
Ilharco, G., Ribeiro, M. T., Wortsman, M., Gururangan, S., Schmidthuber, J., Farhadi, A., and Singer, Y. Editing models with task arithmetic. In \emph{ICLR}, 2023.

\bibitem{yang2024test}
Yang, S., Zhao, H., Jin, X., Yadav, P., and Raffel, C. Test-Time Model Merging on Non-Stationary Streams. In \emph{ICML}, 2024.

\bibitem{yadav2023ties}
Yadav, P., Tam, D., Choset, L., and Raffel, C. TIES-Merging: Resolving Interference When Merging Models. In \emph{NeurIPS}, 2023.

\bibitem{jin2023regmean}
Jin, X., Ren, M., Ye, J., and Han, J. RegMean: Regenerating the Mean for Model Merging. In \emph{ICLR}, 2023.
\end{thebibliography}

\end{document}
